\documentclass[11pt,letterpaper]{article}

\usepackage[margin=1.1in]{geometry}
\usepackage{amsmath, amssymb}
\usepackage{mathtools}
\usepackage{bm}
\usepackage{booktabs}
\usepackage{array}
\usepackage{enumitem}
\usepackage{xcolor}
\usepackage{tcolorbox}
\usepackage{hyperref}
\usepackage[protrusion=true,expansion=false]{microtype}
\usepackage{titlesec}
\usepackage{fancyhdr}
\usepackage{longtable}

\definecolor{accent}{RGB}{30,80,60}
\definecolor{accent2}{RGB}{20,100,100}
\definecolor{boxbg}{RGB}{240,238,232}

\hypersetup{
  colorlinks=true,
  linkcolor=accent,
  citecolor=accent,
  urlcolor=accent2,
  pdftitle={Disc Flight in a Lakefront Internal Boundary Layer (Draft)},
  pdfauthor={charlie}
}

\titleformat{\section}
  {\large\bfseries\color{accent}}
  {\thesection.}{0.6em}{}[\vspace{-2pt}\rule{\linewidth}{0.6pt}]
\titleformat{\subsection}
  {\normalsize\bfseries\color{accent!80!black}}
  {\thesubsection.}{0.5em}{}
\titleformat{\subsubsection}
  {\normalsize\itshape}
  {\thesubsubsection.}{0.4em}{}

\pagestyle{fancy}
\fancyhf{}
\rhead{\small\color{gray}Project 1 $\cdot$ Frisbee + Loyola Wind}
\lhead{\small\color{gray}Manuscript Draft v0.1}
\cfoot{\thepage}
\renewcommand{\headrulewidth}{0.3pt}

\setlength{\parskip}{0.5\baselineskip}
\setlength{\parindent}{0pt}

\title{\bfseries Disc Flight in a Lakefront Internal Boundary Layer:\\Coupling Analytical Frisbee Aerodynamics with Site-Specific Urban CFD\\over Loyola University's Sean Earl Field}
\author{charlie\\\small affiliation TBD}
\date{Draft v0.1 \ $\cdot$\ \today}

\begin{document}
\maketitle

\noindent\textbf{Companion document:} \emph{Project1\_AnalyticalTreatment.tex} (rigorous 6-DOF derivation, stability, gyroscopic precession).

\section*{How to read this document}

Three layers stacked: \textbf{feasibility memo} (\S\ref{sec:feas}), \textbf{project plan} (\S\ref{sec:plan}), and \textbf{paper first draft} (\S\ref{sec:abstract} onward). Read \S\ref{sec:feas} first; if you disagree with the scope changes proposed there, the draft from \S\ref{sec:abstract} onward needs to be reframed.

\section{Feasibility check}
\label{sec:feas}

\subsection{What the notes propose}

\begin{quote}
(1) Analytical treatment of ideal conditions for throwing a frisbee.\\
(2) Build a realistic 3D map of the terrain surrounding Loyola field --- houses, buildings --- simulate wind off the lake spanning 180$^\circ$, generate 3D flow fields across the turf to capture vortices.\\
\emph{Goal:} determine the best conditions for this constrained terrain that best allow for frisbee.\\
\emph{Tools:} OpenFOAM or COMSOL.
\end{quote}

\subsection{Verdict}

\begin{center}
\footnotesize
\begin{tabular}{@{}p{0.32\linewidth}p{0.15\linewidth}p{0.48\linewidth}@{}}
\toprule
\textbf{Component} & \textbf{Verdict} & \textbf{Why} \\
\midrule
Analytical 6-DOF disc model & Solid & Hummel (2003), Crowther \& Potts (2007) give a complete validated 6-DOF formulation with measured coefficients. \\
RANS urban CFD, 13 wind directions $\times$ 15$^\circ$ & Tractable & $\sim$3--10 M cells, 2--8 h per direction on 32 cores, total 1--4 days. Standard COST 732 / AIJ practice; OpenFOAM \texttt{simpleFoam} with k-$\omega$ SST. \\
LES of the same domain & Not realistic for full sweep & $\Delta x \sim 1$--$2$ m near buildings $\Rightarrow$ 50--200 M cells per case. Single direction is feasible (weeks). \\
Coupling CFD wind field into disc trajectory & Novel and tractable & Not done in published literature; the Lukes et al.\ (2022) disc-golf code uses uniform freestream. \emph{This is the novelty hook.} \\
Throwability metric over the field & Tractable as descriptive optimization & Not real-time forecasting; compute mean velocity + TKE at disc altitude, then evaluate trajectory robustness for canonical throws. \\
\bottomrule
\end{tabular}
\end{center}

\subsection{Items in your notes that need adjustment}

\begin{itemize}[leftmargin=*,itemsep=2pt]
  \item \textbf{``Loyola field''} --- your phrasing (``off the lake'') uniquely matches \emph{Sean Earl Field} on the Lake Shore Campus, $\sim$150--250 m from Lake Michigan, not Hoyne Field (a mile inland). \textbf{Confirmed by you: Sean Earl.}
  \item \textbf{``180$^\circ$ span''} --- Sean Earl's lake-facing side is roughly east; a meaningful ``off-lake'' sweep is NNE through ESE ($\sim$90$^\circ$ centered on east). A full 180$^\circ$ is fine if you also characterize the prevailing westerly (off-land) cases for contrast --- think of it as two regimes, not one continuous sweep.
  \item \textbf{``Range finder'' to build the 3D map} --- you don't need a personal LiDAR pass. USGS 3DEP / Illinois Height Modernization Program covers Cook County at 1--8 pts/m$^2$; Microsoft GlobalMLBuildingFootprints + USGS 2023 building-heights release explicitly include Chicago. A handheld rangefinder is useful only for spot-checking the 5--10 buildings nearest the field.
  \item \textbf{``OpenFOAM or COMSOL''} --- pick OpenFOAM. The urban-CFD literature (Blocken, Tominaga, COST 732 community) is overwhelmingly OpenFOAM-based; \texttt{snappyHexMesh} handles dirty STL of campus buildings well; you'll have peer-reviewed cases to validate against.
\end{itemize}

\subsection{The biggest unstated gap}

You do not have a \emph{bridging quantity} that turns a CFD flow field into a throw-quality score. The paper needs one. \S\ref{sec:throwability} proposes a ``throwability metric'' $Q(\mathbf{x}, \theta_{\text{wind}})$ combining mean wind magnitude, turbulence intensity at disc altitude, and predicted lateral trajectory drift for the canonical throws.

\section{Project plan}
\label{sec:plan}

\subsection{Work breakdown ($\approx$14 weeks, single researcher)}

\begin{center}
\footnotesize
\begin{tabular}{@{}llp{0.55\linewidth}@{}}
\toprule
\textbf{Phase} & \textbf{Weeks} & \textbf{Output} \\
\midrule
P0 --- Setup & 1 & OpenFOAM 11 or v2406 installed; Python disc-trajectory simulator stub; QGIS workflow for LiDAR $\to$ STL. \\
P1 --- Analytical disc model & 2--3 & 6-DOF integrator (Python, scipy \texttt{solve\_ivp}). Validated against Hummel (2003) Fig.\ 5-1 throw reconstructions. \\
P2 --- Geometry build & 4 & STL of $\sim$250 m radius around Sean Earl Field, $\sim$5--10 named buildings + tree-canopy proxy. \\
P3 --- RANS CFD sweep & 5--8 & 13 steady-state runs (15$^\circ$ increments). \texttt{simpleFoam} + k-$\omega$ SST + Richards--Hoxey ABL inlet. Output: $U_{\text{mean}}, k$ sampled to 1, 2, 3 m AGL. \\
P4 --- Coupled trajectory & 9--11 & Trilinear-interpolated CFD into 6-DOF disc model. Monte Carlo over launch, heading, spin, release $\alpha$ ($10^4$ throws per direction). \\
P5 --- Synthesis \& manuscript & 12--14 & Throwability map $Q(\mathbf{x}, \theta_{\text{wind}})$; paper revision; figures. \\
\bottomrule
\end{tabular}
\end{center}

\subsection{Compute and storage}

Workstation target: 32 physical cores, 128 GB RAM, $\sim$2 TB scratch SSD. On 16-core / 64 GB expect $\sim$2$\times$ wall-clock and one wind direction per overnight job. Disk budget: $\sim$50 GB per RANS case ($\sim$700 GB for the sweep). LES would multiply by 20--50$\times$.

\subsection{Inputs (all free/accessible)}

\begin{itemize}[leftmargin=*,itemsep=2pt]
  \item \emph{Building footprints/heights:} Cook County GIS, Microsoft GlobalMLBuildingFootprints (ODbL), USGS Overture Maps 2023, OpenStreetMap.
  \item \emph{LiDAR DEM:} USGS 3DEP / Illinois Height Modernization Program (1--8 pts/m$^2$).
  \item \emph{Wind reference:} NWS O'Hare 1991--2020 normals; NOAA GLERL Wilson Avenue Crib (shoreline calibration); Northerly Island ASOS.
  \item \emph{Disc aerodynamic coefficients:} Hummel (2003) Tables 3-1/3-2; Crowther \& Potts (2007).
\end{itemize}

\subsection{Risks \& mitigations}

\begin{center}
\footnotesize
\begin{tabular}{@{}p{0.46\linewidth}p{0.46\linewidth}@{}}
\toprule
\textbf{Risk} & \textbf{Mitigation} \\
\midrule
k-$\omega$ SST under-predicts rooftop separation at sharp roof corners (well-known). & Validate one direction against a LES spot check; follow Blocken (2018) on RANS limits. Don't claim peak/fluctuating quantities --- report means + TKE only. \\
Disc-coefficient uncertainty at off-design $\alpha$ ($>20^\circ$). & Restrict throw library to $\alpha \in [-5^\circ, +20^\circ]$; flag hammer (high-$\alpha$) results as exploratory. \\
Horizontal inhomogeneity of ABL inlet (Hargreaves \& Wright 2007). & Use Yang et al.\ (2009) corrected wall functions; verify inlet profile preserved upstream of buildings. \\
Sean Earl is not a competition-grade Ultimate field. & Reframe as a methods paper using Sean Earl as exemplar; methodology generalizes. \\
Reviewers ask for measurement validation. & Use Tominaga (2008) AIJ-B cube benchmark for the CFD; cite Hummel's marker-tracked flights for the disc model. Both are public. \\
\bottomrule
\end{tabular}
\end{center}

\section{Abstract (draft)}
\label{sec:abstract}

Outdoor disc sports are dominated by wind, yet the published aerodynamic and trajectory literature for flying discs (Hummel 2003; Potts \& Crowther 2002; Lukes et al.\ 2022) uses uniform-freestream assumptions that ignore site-specific flow structure. Conversely, the urban-CFD literature (Tominaga et al.\ 2008; Blocken 2014, 2018) characterizes site-specific wind fields but stops at pedestrian-level mean quantities. Here we couple a validated 6-degree-of-freedom (6-DOF) disc-flight model to steady-RANS atmospheric CFD of a lakefront university campus (Loyola University Chicago, Sean Earl Field; $\sim$200 m from the Lake Michigan shore). We sweep 13 inflow directions in 15$^\circ$ increments using OpenFOAM with k-$\omega$ SST turbulence and a Richards--Hoxey atmospheric-boundary-layer inlet, validated against COST Action 732 benchmarks. Disc trajectories for canonical throws (backhand, forehand, hammer; 175~g Ultrastar) are integrated through the resulting time-averaged + turbulence-intensity fields. We report a spatial ``throwability'' map $Q(\mathbf{x}, \theta_{\text{wind}})$ and identify the lakefront internal boundary layer --- formed by the sharp roughness transition at the shoreline --- as the dominant local feature controlling disc-altitude flow. The methodology generalizes to any outdoor sport played on a roughness-discontinuity site.

\section{Introduction (draft)}

The Frisbee is one of the most-studied non-airfoil aerodynamic objects in sports physics. Stilley \& Carstens (1972) provided the first parametric model; Yasada (1999), Hummel (2003), and Potts \& Crowther (2002) supplied successively more complete wind-tunnel measurements of the lift, drag, and pitching-moment coefficients, and Hummel's UC Davis thesis remains the most-cited 6-DOF flight reconstruction. Subsequent CFD validation by Kamaruddin (2011) and Lukes et al.\ (2022, 2023) using RANS k-$\omega$ SST in OpenFOAM has converged the modeling stack to within a few percent of wind-tunnel $C_L$, $C_D$, and $C_M$ at moderate angles of attack.

Yet every published disc-trajectory study assumes a spatially uniform, steady freestream. As Lukes et al.\ (2022) explicitly note, ``the biggest uncertainty in trajectory experiments is the wind environment.'' This uncertainty is not noise to be averaged away: outdoor disc sports --- Ultimate, disc golf --- are played in geographically structured turbulent flow. A throw whose intended landing zone lies in a sheltered eddy behind a building experiences a categorically different flow than the same throw 20 m to the side.

Independently, the urban-CFD community has developed mature best-practice guidance for site-specific pedestrian-level wind: the COST Action 732 guidebook (Franke et al.\ 2007), the AIJ guidelines (Tominaga et al.\ 2008), and the 50-year retrospective of Blocken (2014). Standard practice uses steady RANS with k-$\omega$ SST or realizable k-$\varepsilon$ on \texttt{snappyHexMesh}-generated meshes of $\sim 10^7$ cells, with Richards--Hoxey atmospheric inlets. This work has shaped pedestrian-comfort regulation in dozens of cities, but to our knowledge has never been applied to outdoor sports trajectory prediction.

The present work fills that gap. We choose Sean Earl Field at Loyola University Chicago's Lake Shore Campus as a study site because it offers an unusually clean physical setup: the field sits within $\sim$250 m of Lake Michigan, with no major intervening built obstacle to the east, behind it a typical low-rise urban-campus block to the west. Easterly winds therefore arrive after $\sim$80 km of unbroken over-water fetch ($z_0 \approx 0.001$ m), pass through a sharp roughness transition at the shoreline, and develop a shallow internal boundary layer over the field. Westerly winds arrive through the campus block. The site cleanly separates two regimes (smooth-fetch and urban-fetch) within a 180$^\circ$ heading sweep.

\textbf{Our contributions:}
\begin{enumerate}[leftmargin=*,itemsep=2pt]
  \item The first published two-way coupling of an environmental-CFD wind field to a validated 6-DOF disc-flight model.
  \item A ``throwability metric'' $Q(\mathbf{x}, \theta_{\text{wind}})$ that compresses the spatial CFD field plus trajectory uncertainty into an actionable single scalar per (launch position, wind heading).
  \item A case study identifying the lakefront internal boundary layer as the dominant local control on Ultimate-disc flight at the chosen site.
\end{enumerate}

\section{Analytical disc aerodynamics (summary)}

Full derivation in the companion document. The disc carries a 6-DOF state $(\mathbf{x}, \mathbf{v}, \psi, \theta, \phi_s, p, q, \Omega)$ governed by Newton's equation (\ref{eq:trans-summary}) plus Euler's equation for a symmetric rotor (\ref{eq:rot-summary}). In the high-spin limit relevant to all real throws, the transverse angular rates are slaved algebraically to the applied aerodynamic moments,
\begin{equation}
p = -\frac{M}{I_p \Omega}, \qquad q = \frac{R}{I_p \Omega},
\end{equation}
recovering the textbook gyroscopic-precession picture. The pitching moment vanishes at the trim angle
\begin{equation}
\alpha_{\text{trim}} = \alpha_0 - \frac{C_{M0}}{C_{M\alpha}} \approx 6.7^\circ \text{ (Hummel coeffs)},
\end{equation}
which is the asymptotic glide angle of attack; the corresponding glide $L/D \approx 3.4$. The full translational and rotational equations entering the numerical simulator are
\begin{align}
m\,\dot{\mathbf{v}} &= -mg\,\hat{\mathbf{n}}_3 + \tfrac{1}{2}\rho V^2 A\big[\,C_L(\alpha)\hat{\mathbf{L}} - C_D(\alpha)\hat{\mathbf{v}}_{\text{rel}}\big], \label{eq:trans-summary}\\
I_p \dot{p} + I_p q \Omega &= R,\qquad I_p \dot{q} - I_p p \Omega = M,\qquad I_s \dot{\Omega} = N. \label{eq:rot-summary}
\end{align}

\section{Site description: Sean Earl Field}

Sean Earl Field is a small multi-use turf flanked by Mertz Residence Hall (W) and the Norville Center for Intercollegiate Athletics (S), with a clear lakefront promenade and beach to the east. Per the Loyola Lake Shore Campus map (LUC 2025), the field's eastern sideline lies approximately 150--250 m from the Lake Michigan shoreline; the field's long axis runs roughly N--S. Loyola explicitly lists Ultimate among the field's sanctioned uses.

\paragraph{Climatology.} Chicago lies in the mid-latitude westerly belt; the annual prevailing direction at ORD is WSW--SW (Illinois State Water Survey wind roses). The 1991--2020 NWS annual mean speed at O'Hare is $\sim$9--10 mph. However, ORD is 27 km inland; lakefront stations (NOAA GLERL Wilson Avenue Crib, Northerly Island ASOS) record higher mean speed and gustier easterlies, especially during the warm-season lake-breeze regime.

\paragraph{Why this site.} Easterly inflow at Sean Earl Field passes over $\sim$80 km of unbroken Lake Michigan fetch (to St.\ Joseph / Benton Harbor MI), giving a near-neutral, very-low-$z_0$ marine boundary layer. The shoreline imposes a sharp $z_0$ transition ($\sim 10^{-3}$ to $\sim$0.3 m). The downwind internal boundary layer thickness $\delta_{\text{IBL}}$ grows as $\delta_{\text{IBL}} \approx 0.10\,x^{0.8}$ (Garratt 1990), reaching $\sim$5 m by the eastern sideline and $\sim$7--10 m at midfield --- fully encompassing the 1--3 m disc-altitude band but in a sheared, non-equilibrium state. \emph{This is the dominant local physical feature.}

\section{CFD methodology}

\paragraph{Geometry.} A $\sim$500$\times$500$\times$80 m domain centered on Sean Earl Field, built from Microsoft GlobalMLBuildingFootprints + USGS 2023 heights for the campus block and 2--3 nearest residential blocks west of campus, with USGS 3DEP LiDAR DEM for ground. Buildings imported as STL after height-extrusion; tree canopy on the lakefront strip modeled as a porous zone with leaf area density $\alpha_{\text{LAD}} \approx 0.5$ m$^2$/m$^3$ (Gromke \& Ruck 2008). Domain extent satisfies COST 732 / Tominaga (2008) guidelines: 5H upstream/lateral/top, 15H downstream of tallest building ($H \approx 22$ m for Mertz). Blockage $<3\%$.

\paragraph{Mesh.} \texttt{snappyHexMesh} hex-dominant. Background cell 4 m, three refinement levels in a 250 m field-centered box (0.5 m near the field surface, 1 m at building faces, surface layers giving $y^+ \approx 30$--$100$ for wall functions). Total $\approx$ 5--8 M cells.

\paragraph{Solver.} OpenFOAM v2406 \texttt{simpleFoam}, k-$\omega$ SST, second-order schemes. Richards \& Hoxey (1993) ABL inlet with Yang et al.\ (2009) corrected wall functions to suppress horizontal inhomogeneity. Inlet roughness: $z_0 = 0.001$ m (water) for headings in $[60^\circ, 120^\circ]$ east sector; $z_0 = 0.5$ m (urban) for $[180^\circ, 360^\circ]$; transitional zones blended. Top BC constant-shear; sides symmetry; outlet pressure-outlet. Reference: $u_{\text{ref}} = 5$ m/s at $z_{\text{ref}} = 10$ m (representative shoreline mean).

\paragraph{Cases.} Inflow heading $\theta_{\text{wind}} \in \{0^\circ, 15^\circ, 30^\circ, \ldots, 180^\circ\}$ (13 cases). Steady RANS to residuals $<10^{-4}$. Wall-clock per case $\approx 4$ h on 32 cores.

\paragraph{Validation.} (i) ABL profile recovery at empty-domain upstream of buildings (Hargreaves \& Wright 2007 protocol). (ii) Tominaga (2008) AIJ-B cube benchmark comparison for the \texttt{snappyHexMesh} + k-$\omega$ SST stack on a separate test mesh. (iii) Loyola-site spot checks: a handheld anemometer at three field-corner locations during 2026 summer/fall lake-breeze events, matched against the corresponding CFD case interpolated to instrument height.

\section{Coupled trajectory simulation}

For each $\theta_{\text{wind}}$, the CFD output is trilinearly interpolated onto a $0.5\,\text{m}^3$ grid covering 1--6 m AGL across the field. The 6-DOF integrator evaluates $\mathbf{U}(\mathbf{x})$ at the disc's instantaneous position; the steady-RANS field provides $U_{\text{mean}}$, and turbulence intensity $I = \sqrt{2k/3}/U_{\text{mean}}$ is mapped into a stochastic gust contribution $\mathbf{u}'(t)$ generated by a Mann (1998) synthetic-turbulence Fourier sampler with the local $I$ and integral length scale set by $k/\varepsilon$ from the RANS.

For each $(\theta_{\text{wind}}, \text{throw type})$, $10^4$ Monte Carlo throws are launched from a 5$\times$5 m grid of launch points spanning the field with realistic release-parameter scatter:
\begin{itemize}[leftmargin=*,itemsep=1pt]
  \item Speed $V_0 \sim \mathcal{N}(15, 1.5)$ m/s for backhand and forehand; $\mathcal{N}(11, 1.5)$ for hammer.
  \item Spin $\Omega_0 \sim \mathcal{N}(100, 10)$ rad/s.
  \item Release pitch $\alpha_0 \sim \mathcal{N}(+8^\circ, 4^\circ)$ backhand; $\mathcal{N}(+6^\circ, 4^\circ)$ forehand; $\mathcal{N}(+40^\circ, 8^\circ)$ hammer.
  \item Heading bearing $\theta_{\text{throw}}$ varied over 0--360$^\circ$ in 30$^\circ$ bins.
\end{itemize}
Landing-point distributions are characterized by 90\%-containment ellipse parameters and a ``delivered-to-target'' success rate (landing within 3 m of intended target at canonical 25 m throws).

\section{Throwability metric (the novelty)}
\label{sec:throwability}

For each launch position $\mathbf{x}$ and each wind heading $\theta_{\text{wind}}$:
\begin{equation}
Q(\mathbf{x}, \theta_{\text{wind}}) =
\underbrace{\Big(1 - \frac{\langle |U_{2\text{m}}| \rangle_{\text{field}}}{U_{\text{ref}}}\Big)}_{\text{wind-magnitude penalty}}
\cdot
\underbrace{\big(1 - I_{2\text{m}}(\mathbf{x})\big)}_{\text{turbulence penalty}}
\cdot
\underbrace{\overline{S}_{\text{throws}}(\mathbf{x}, \theta_{\text{wind}})}_{\text{simulated success rate}}
\end{equation}
with each factor normalized to $[0, 1]$. $Q$ ranges from 0 (unplayable) to 1 (calm-air-equivalent). Maps of $Q$ over the field surface, parameterized by $\theta_{\text{wind}}$, form the principal result.

\section{Expected results / placeholder findings}

To be populated after the simulation sweep. Hypothesized:
\begin{itemize}[leftmargin=*,itemsep=2pt]
  \item Easterly cases ($90^\circ \pm 30^\circ$) show a sharp $Q$ gradient across the field's eastern half corresponding to the developing IBL; backhand throws from the windward sideline are most consistent; hammer throws systematically degraded.
  \item Westerly cases ($270^\circ \pm 30^\circ$) show patchier $Q$ with sheltered ``throw lanes'' downstream of Norville Center; forehand throws benefit most.
  \item Best-condition wind direction for canonical 25 m throws is hypothesized to be a light SSW ($\sim$210$^\circ$, 3--5 m/s reference), aligning with prevailing climatology --- a satisfying convergence of meteorology and playability.
\end{itemize}

\section{Discussion (to draft after results)}

Generalizability to other outdoor sports; steady-RANS limitations (peak gusts and spatially correlated gust structure not captured); LES of one or two critical directions is the obvious follow-on; coaching implications.

\section{Open questions for charlie}

\begin{enumerate}[leftmargin=*,itemsep=2pt]
  \item One LES validation run for one wind direction --- worth $\sim$2 weeks for a much stronger paper?
  \item Are you committed to the 2026 summer-anemometer measurement plan? Realistic field-access logistics?
  \item Co-authors? Department affiliation? Funding to acknowledge?
  \item Hammer (high-$\alpha$) is the throw most affected by wind \emph{and} the one whose aerodynamic coefficients are least well-validated. OK to flag as exploratory, or drop?
\end{enumerate}

\section*{Key references (verified)}

\begin{description}[font=\normalfont\itshape,leftmargin=1em]
  \item Blocken, B.\ (2014). 50 years of computational wind engineering: Past, present and future. \emph{JWEIA} \textbf{129}, 69--102.
  \item Blocken, B.\ (2018). LES over RANS in building simulation\ldots\ \emph{Building Simulation} \textbf{11}, 821--870.
  \item Crowther, W.J.\ \& Potts, J.R.\ (2007). Simulation of a spin-stabilised sports disc. \emph{Sports Engineering} \textbf{10}(1), 3--21.
  \item Franke, J.\ et al.\ (2007). \emph{Best Practice Guideline for CFD Simulation of Flows in the Urban Environment.} COST Action 732.
  \item Hargreaves, D.M.\ \& Wright, N.G.\ (2007). On the use of the k-$\varepsilon$ model in commercial CFD software\ldots\ \emph{JWEIA} \textbf{95}, 355--369.
  \item Hubbard, M.\ \& Hummel, S.A.\ (2000). Simulation of Frisbee Flight. 5th Math \& Computers in Sport, Sydney.
  \item Hummel, S.A.\ (2003). \emph{Frisbee Flight Simulation and Throw Biomechanics.} MS Thesis, UC Davis.
  \item Kamaruddin, N.M., Potts, J.R.\ \& Crowther, W.J.\ (2018). Aerodynamic performance of flying discs. \emph{Aircraft Eng.\ Aerosp.\ Tech.}\ \textbf{90}(2), 390--397.
  \item Lissaman, P.B.S.\ \& Hubbard, M.\ (2010). Maximum range of flying discs. \emph{Procedia Engineering} \textbf{2}(2), 2529--2535.
  \item Lukes, R., Hart, J.\ et al.\ (2022). Disc golf trajectory modelling combining CFD and rigid body dynamics. \emph{Sports Engineering} \textbf{25}, 24.
  \item Potts, J.R.\ \& Crowther, W.J.\ (2002). Frisbee\texttrademark{} aerodynamics. AIAA 2002-3150.
  \item Richards, P.J.\ \& Hoxey, R.P.\ (1993). Appropriate boundary conditions for computational wind engineering models. \emph{JWEIA} \textbf{46--47}, 145--153.
  \item Tominaga, Y.\ et al.\ (2008). AIJ guidelines for practical applications of CFD to pedestrian wind environment. \emph{JWEIA} \textbf{96}, 1749--1761.
  \item Yang, Y.\ et al.\ (2009). New inflow boundary conditions for modelling the neutral equilibrium atmospheric boundary layer. \emph{JWEIA} \textbf{97}, 88--95.
\end{description}

\end{document}
ment}
